\fichatitulo{20 · RECUPERAÇÃO}{Livro · edição on-line 2009}{Avaliação em recuperação da informação}
{\small Manning; Raghavan; Schütze. \textit{Introduction to Information Retrieval}. Livro · edição on-line 2009. \href{https://nlp.stanford.edu/IR-book/pdf/irbookonlinereading.pdf}{Fonte consultada}.\par}

\campo{Problema e objetivo}
Como avaliar a relevância dos documentos recuperados?

\campo{Principal contribuição · síntese interpretativa}
Sistematiza coleções de teste, necessidades de informação, julgamentos de relevância e medidas de recuperação, distinguindo-as da utilidade para o usuário.

\campo{Método / natureza da evidência}
Exposição didática de conceitos e procedimentos; recorte de capítulo, não fichamento integral do livro.

\campo{Resultado ou argumento principal}
A avaliação exige coleção, consultas vinculadas a necessidades e julgamentos. Medidas de conjuntos diferem de medidas que consideram a posição dos resultados.

\campo{Excertos-chave · seleção literal}
\begin{itemize}
\excerto{1}{Relevance is assessed relative to an information need, not a query.}{Seção 8.1.}
\excerto{2}{These information needs are best designed by domain experts.}{Seção 8.5.}
\end{itemize}

\campo{Limitações e análise crítica}
Autores: julgar relevância tem custo, e coleções grandes usam avaliações parciais. Análise da ficha: resultados dependem da representatividade das perguntas e do protocolo de anotação.

\campo{Aproveitamento na dissertação · proposta}
Metodologia: definir unidade recuperada, relevância, profundidade de avaliação e avaliação humana.

\registro{Fonte consultada nesta preparação: Capítulo 8: trechos das seções 8.1, 8.3, 8.4, 8.5 e 8.6. Livro originalmente publicado em 2008; PDF on-line identificado como 2009. Verificação estrutural automática indisponível; usados localizadores textuais por seção. Chave: \texttt{manningEtAl2008informationRetrieval}.}
